\begin{definition}[Длина границы грани]
Длина границы — число сторон рёбер при обходе грани с учётом кратности.
Мост учитывается дважды на границе одной грани. Поэтому
\(\sum_F \ell(F)=2m\): у каждого ребра две стороны.
\end{definition}
\begin{graphillustration}{Пример обхода границ: мост $cv$ встречается на границе внешней грани дважды — по одной стороне в каждом направлении.}{fig:15-boundaries}
\begin{tikzpicture}[font=\small]
\node[gi vertex] (a) at (0,0) {$a$};
\node[gi vertex] (b) at (2,0) {$b$};
\node[gi vertex] (c) at (1,1.7) {$c$};
\node[gi vertex] (v) at (1,3) {$v$};
\draw[gi edge] (a) -- (b);
\draw[gi edge] (b) -- (c);
\draw[gi edge] (c) -- (a);
\draw[gi edge] (c) -- (v);
\draw[gi flow] (.8,2)--(.8,2.65);\draw[gi flow] (1.2,2.65)--(1.2,2);
\node[gi note] at (1,0.65) {$F_1$};
\node[gi note] at (3,1.8) {$F_0$};
\node[gi note] at (6.3,1.2) {$\ell(F_1)=3$\\$\ell(F_0)=3+2=5$\\[5pt]$\ell(F_0)+\ell(F_1)=8=2m$};
\end{tikzpicture}
\end{graphillustration}


\begin{lemma}[Оценки числа рёбер]
Для простого планарного графа с \(n\ge3\)
\[
 m\le3n-6.
\]
Если граф двудолен, то \(m\le2n-4\).
\end{lemma}
\begin{proof}
Пусть сначала граф связен. Если это дерево, \(m=n-1\), и обе оценки
следуют из \(n\ge3\). Иначе любая граница грани имеет длину не менее 3:
длина 1 требовала бы петли, длина 2 — кратных рёбер или единственного
ребра, дважды обходимого как весь связный граф. Последнее невозможно
при \(n\ge3\). В двудольном графе все замкнутые маршруты имеют чётную
длину, так что длина грани не менее 4.

В первом случае \(3f\le2m\), и по Эйлеру
\(2=n-m+f\le n-m/3\), откуда \(m\le3n-6\).
Во втором \(4f\le2m\), следовательно,
\(2\le n-m/2\), что даёт \(m\le2n-4\).
Если граф несвязен, соединим его компоненты непересекающимися мостами.
Простота сохраняется; в двудольном случае при необходимости меняем
названия долей в одной из соединяемых компонент. Получим связный
простой (соответственно двудольный) планарный граф на тех же вершинах
с не меньшим числом рёбер. Его оценка влечёт исходную.
\end{proof}
\begin{theorem}
Графы \(K_5\) и \(K_{3,3}\) непланарны.
\end{theorem}
\begin{proof}
У \(K_5\) имеем \(n=5,m=10\), но для планарного графа
\(m\le3\cdot5-6=9\), противоречие.
У двудольного \(K_{3,3}\) имеем \(n=6,m=9\), но
\(m\le2\cdot6-4=8\), снова противоречие.
\end{proof}
\begin{center}
\begin{tikzpicture}[every node/.style={circle,fill=cbBlue,inner sep=2pt}]
\foreach \i in {1,...,5}{\node (p\i) at ({90+72*(\i-1)}:1.15) {};}
\foreach \i in {1,...,5}{\foreach \j in {1,...,5}{\ifnum\i<\j\draw (p\i)--(p\j);\fi}}
\node[fill=none,rectangle] at (0,-1.6) {$K_5$};
\foreach \i in {1,2,3}{\node (a\i) at (3,{2-\i}) {};\node (b\i) at (5,{2-\i}) {};}
\foreach \i in {1,2,3}{\foreach \j in {1,2,3}{\draw (a\i)--(b\j);}}
\node[fill=none,rectangle] at (4,-1.6) {$K_{3,3}$};
\end{tikzpicture}
\end{center}
Пересечения на этих рисунках сами по себе не доказывают непланарность:
доказательство должно исключать \emph{любую} укладку, что и делают оценки.
\begin{cor}
Ни один планарный граф не содержит подграфа, являющегося подразбиением
\(K_5\) или \(K_{3,3}\).
\end{cor}
\begin{proof}
Подграф планарного графа планарен: достаточно удалить лишние дуги и точки
из укладки. Планарность подразбиения равносильна планарности исходного
графа (билет 13), тогда \(K_5\) или \(K_{3,3}\) был бы планарен.
\end{proof}
